\batchmode
\makeatletter
\def\input@path{{"C:/Users/tkpme/Dropbox/BookPCRA/Ch 0 Preface/"}}
\makeatother
\documentclass[graybox,envcountchap,sectrefs]{svmono}\usepackage[]{graphicx}\usepackage[]{xcolor}
% maxwidth is the original width if it is less than linewidth
% otherwise use linewidth (to make sure the graphics do not exceed the margin)
\makeatletter
\def\maxwidth{ %
  \ifdim\Gin@nat@width>\linewidth
    \linewidth
  \else
    \Gin@nat@width
  \fi
}
\makeatother

\definecolor{fgcolor}{rgb}{0.345, 0.345, 0.345}
\newcommand{\hlnum}[1]{\textcolor[rgb]{0.686,0.059,0.569}{#1}}%
\newcommand{\hlsng}[1]{\textcolor[rgb]{0.192,0.494,0.8}{#1}}%
\newcommand{\hlcom}[1]{\textcolor[rgb]{0.678,0.584,0.686}{\textit{#1}}}%
\newcommand{\hlopt}[1]{\textcolor[rgb]{0,0,0}{#1}}%
\newcommand{\hldef}[1]{\textcolor[rgb]{0.345,0.345,0.345}{#1}}%
\newcommand{\hlkwa}[1]{\textcolor[rgb]{0.161,0.373,0.58}{\textbf{#1}}}%
\newcommand{\hlkwb}[1]{\textcolor[rgb]{0.69,0.353,0.396}{#1}}%
\newcommand{\hlkwc}[1]{\textcolor[rgb]{0.333,0.667,0.333}{#1}}%
\newcommand{\hlkwd}[1]{\textcolor[rgb]{0.737,0.353,0.396}{\textbf{#1}}}%
\let\hlipl\hlkwb

\usepackage{framed}
\makeatletter
\newenvironment{kframe}{%
 \def\at@end@of@kframe{}%
 \ifinner\ifhmode%
  \def\at@end@of@kframe{\end{minipage}}%
  \begin{minipage}{\columnwidth}%
 \fi\fi%
 \def\FrameCommand##1{\hskip\@totalleftmargin \hskip-\fboxsep
 \colorbox{shadecolor}{##1}\hskip-\fboxsep
     % There is no \\@totalrightmargin, so:
     \hskip-\linewidth \hskip-\@totalleftmargin \hskip\columnwidth}%
 \MakeFramed {\advance\hsize-\width
   \@totalleftmargin\z@ \linewidth\hsize
   \@setminipage}}%
 {\par\unskip\endMakeFramed%
 \at@end@of@kframe}
\makeatother

\definecolor{shadecolor}{rgb}{.97, .97, .97}
\definecolor{messagecolor}{rgb}{0, 0, 0}
\definecolor{warningcolor}{rgb}{1, 0, 1}
\definecolor{errorcolor}{rgb}{1, 0, 0}
\newenvironment{knitrout}{}{} % an empty environment to be redefined in TeX

\usepackage{alltt}
\usepackage[T1]{fontenc}
\usepackage[utf8]{inputenc}
\usepackage{geometry}
\geometry{verbose,tmargin=1in,bmargin=1in,lmargin=1in,rmargin=1in}
\setcounter{tocdepth}{3}
\setlength{\parskip}{\medskipamount}
\setlength{\parindent}{0pt}
\usepackage{float}
\usepackage{textcomp}
\usepackage{url}
\usepackage{setspace}
\onehalfspacing
\usepackage[unicode=true,pdfusetitle,
 bookmarks=true,bookmarksnumbered=true,bookmarksopen=false,
 breaklinks=false,pdfborder={0 0 1},backref=false,colorlinks=false]
 {hyperref}

\makeatletter
%%%%%%%%%%%%%%%%%%%%%%%%%%%%%% User specified LaTeX commands.
%%%%%%%%%%%%%%%%%%%% book.tex %%%%%%%%%%%%%%%%%%%%%%%%%%%%%
%
% sample root file for the chapters of your "monograph"
%
% Use this file as a template for your own input.
%
%%%%%%%%%%%%%%%% Springer-Verlag %%%%%%%%%%%%%%%%%%%%%%%%%%


% RECOMMENDED %%%%%%%%%%%%%%%%%%%%%%%%%%%%%%%%%%%%%%%%%%%%%%%%%%%


% choose options for [] as required from the list
% in the Reference Guide


\usepackage[bottom]{footmisc}% places footnotes at page bottom

\usepackage{hyperref}
\def\UrlBreaks{\do\/\do-}

%\usepackage[ style=authoryear,natbib=true,giveninits=true,backend=biber,maxcitenames=2]{biblatex}
\usepackage[style=authoryear, natbib=true, firstinits=false, backend=biber, maxcitenames=2]{biblatex}


\setlength{\bibitemsep}{0.5\baselineskip}
\DeclareNameAlias{sortname}{family-given}

\IfFileExists{C:/Users/Doug/Dropbox/BookPCRA/mpslbookCombined.bib}
{\addbibresource[location=local]{C:/Users/Doug/Dropbox/BookPCRA/mpslbookCombined.bib}}
{   \IfFileExists{C:/Users/tkpme/Dropbox/BookPCRA/mpslbookCombined.bib}
    {\addbibresource[location=local]{C:/Users/tkpme/Dropbox/BookPCRA/mpslbookCombined.bib}}
    {   \IfFileExists{C:/Users/stoya/Dropbox/BookPCRA/mpslbookCombined.bib}
        {\addbibresource[location=local]{C:/Users/stoya/Dropbox/BookPCRA/mpslbookCombined.bib}}
        {   \IfFileExists{C:/Users/4thauthor/Dropbox/BookPCRA/mpslbookCombined.bib}
            {\addbibresource[location=local]{C:/Users/4thauthor/Dropbox/BookPCRA/mpslbookCombined.bib}}
            {}
        }
    }
}
\renewcommand\cite{\citet}

\usepackage{amsmath}
\usepackage{amssymb}
\usepackage[p,osf]{stickstootext}
\usepackage[stix2,vvarbb]{newtxmath}
\usepackage[varqu,varl]{inconsolata}
% \usepackage{txfonts}
\usepackage{upgreek}

\DeclareMathOperator{\argmax}{arg\,max}
\DeclareMathOperator{\argmin}{arg\,min}
\usepackage{array}
\usepackage{booktabs}
\usepackage{threeparttablex}

\usepackage[utf8]{inputenc}
%%%%%%%%%%%%%%%%%%%%%%%%%%%%%%%%%%%%%%%%%%%
\IfFileExists{C:/Users/Doug/Dropbox/BookPCRA/macroSP_August_05_2021.tex}
{\input{C:/Users/Doug/Dropbox/BookPCRA/macroSP_August_05_2021.tex}}
{   \IfFileExists{C:/Users/tkpme/Dropbox/BookPCRA/macroSP_August_05_2021.tex}
    {\input{C:/Users/tkpme/Dropbox/BookPCRA/macroSP_August_05_2021.tex}}
    {   \IfFileExists{C:/Users/stoya/Dropbox/BookPCRA/macroSP_August_05_2021.tex}
        {\input{C:/Users/stoya/Dropbox/BookPCRA/macroSP_August_05_2021.tex}}
        {   \IfFileExists{C:/Users/4thAuthor/Dropbox/BookPCRA/macroSP_August_05_2021.tex}
            {\input{C:/Users/4thAuthor/Dropbox/BookPCRA/macroSP_August_05_2021.tex}}
            {}
        }
    }
}

%\belowsink=90pt

\hypersetup{pdfstartview={XYZ null null 1.25}}


%%%%%%%%%%%%%%%%%%%%%%%%%%%%%%%%%%%%%%%%%%%%%

\makeatother
\IfFileExists{upquote.sty}{\usepackage{upquote}}{}
\begin{document}
\title{Robust Portfolio Construction and Risk Analysis}
\author{R. Douglas Martin, Thomas K. Philips, Bernd Scherer and Kirk Li}
\subtitle{Preface}

\maketitle
\newpage{}

\addcontentsline{toc}{chapter}{PREFACE}

This Robust Portfolio Construction and Risk Analysis (RPCRA) book,
with its unique \emph{reproducibility} code, \emph{robust estimators},
and \emph{special focus} topics, has been written for use in advanced
upper level undergraduate, M.S., and PhD level quantitative finance
courses. It grew out of our collective experiences as researchers,
educators and investors, and is thus supportive of research projects
in such programs, as well as in commercial quantitative portfolio
management and research.

Much of the content in Chapters 1-9 is suitable for a full academic
year M.S. course, and selected portions of these chapters can serve
the needs of a one-semester course. The additional material in Chapters
10-16 is well suited to special topics and project oriented courses.
The chapters and the topics they cover are:
\begin{description}
\item [{Chapter}] 1. Introduction to Portfolio Management
\item [{Chapter}] 2. Foundations
\item [{Chapter}] 3. Modeling Asset Returns
\item [{Chapter}] 4. Covariance Matrices
\item [{Chapter}] 5. Portfolio Constraints and Penalties
\item [{Chapter}] 6. Time Series Models
\item [{Chapter}] 7. Cross Section Factor Models
\item [{Chapter}] 8. Portfolio Risk Analysis
\item [{Chapter}] 9. Minimum Downside Risk Portfolios
\item [{Chapter}] 10. Performance Evaluation
\item [{Chapter}] 11. Expected Returns
\item [{Chapter}] 12. Fixed Income Portfolio Management
\item [{Chapter}] 13. Portfolios With Options
\item [{Chapter}] 14. Managing Portfolios to Meet Spending Objectives
\item [{Chapter}] 15. Influence Functions and Standard Errors
\item [{Chapter}] 16. Robust Estimators
\end{description}
In addition, there are five appendices
\begin{description}
\item [{Appendix}] A. Mathematical Tools For PCRA
\item [{Appendix}] B. Mathematical Tools For PCRA
\item [{Appendix}] C. Utility Theory
\item [{Appendix}] D. Portfolio Theory
\item [{Appendix}] E. Machine Learning
\end{description}
In addition, we will maintain a companion website for this book at
\url{https://www.robustpcra.org}. On this site, we will post articles,
provide code and examples, and point to new uses of the datasets included
in the \texttt{PCRA} package.

\section*{Reproducibility}

R packages that include R functions and data support reproducibility
of figures and tables throughout the book. Specifically, this is the
case for most of the figures and tables in Chapters 2-6 and 9-12,
as well as the figures and tables in some sections of Chapters 1 and
7-8. These features allow students to get hands-on experience in replicating
the figures and tables in RPCRA, and in extending the code to study
portfolio performance for different different time periods, different
portfolio stock groups (such as capitalization groups and sectors),
and different returns frequencies. These features also support instructors
in creating new computing-oriented portfolio construction and risk
analysis exercises.

\subsection*{R Packages for the RPCRA Book and its Reproducibility Character}

The following R packages were developed specifically for the writing
of the RPCRA book and its reproducibility character, and are downloadable
from the CRAN website \url{https://cran.r-project.org}:
\begin{itemize}
\item \texttt{PCRA}:\quad{}Portfolio construction and risk analysis
\item \texttt{facmodCS}:\quad{}Cross-section factor models
\item \texttt{facmodTS}:\quad{}Time series factor models
\item \texttt{robustGarch}:\quad{}Robust Garch(1,\ 1) models.
\end{itemize}

\subsection*{Data for Reproducibility}

The \texttt{PCRA} package contains data sets that support reproducibility
as well as a multitude of exercises at the end of each chapter. They
may also be used by course intructors to create new exercises. The
datasets include:
\begin{itemize}
\item Center for Research of Security Prices (\CRSPTM) monthly, weekly
and daily stocks data for 294 stocks from 1993 to 2015, and Standard
and Poors Global Market Intelligence (SPGMI) monthly factors data
for the \CRSPTM stocks data (for reproducibility of results in Chapters
2-9)
\item S\&P, CBOE, and Hedge Fund data sets for reproducibility in Chapters
10 and 11.
\item ICE Fixed Income datasets for reproducibility in Chapter 12
\end{itemize}
An overview of the \texttt{PCRA} package and its data sets are provided
in the \emph{PCRA Package and Data Overview} vignette that is downloadable
from \url{https://cran.r-project.org/web/packages/PCRA}. Detailed
information about the \CRSPTM\ monthly, weekly and daily stocks
data, and \texttt{factorsSPGMI} monthly data is provided in the \emph{CRSP
Stocks and SPGMI Factors in PCRA} vignette that is also downloadable
from the above \texttt{PCRA} package link.

\section*{Robust Estimators}

Having firsthand experience of the impact of outliers in portfolio
management, we were struck by how little attention has been paid to
this important problem in portfolio management textbooks, and more
generally in quantitative finance books. More often than not, students
learn the elegant mathematics of portfolio theory with little appreciation
for its fragilityin the face of the prodigious levels of noise that
permeate most markets. Such noise includes outliers that can have
large distorting influences on the classical estimators that are routinely
used in risk analysis.

The adverse influence of outliers on these classical estimators can
be avoided by using modern robust alternatives that are not much influenced
by outliers, and which provide better fits to the bulk of the data
than do classical estimators. Chapters 2, 4, 6, and 7 introduce modern
robust estimators of returns means and standard errors, covariance
matrices, time-series factor models, and cross-section factor models
respectively. These chapters provide clear examples of the superior
fits afforded by robust estimators relative to classical estimators
in the presence ofeven a few outliers. Chapter 16 provides background
mathematical details on robust estimators and provides a bridge to
RobStatTM, the \citet{Maronna2019} \emph{Robust Statistics: Theory
and Methods} book that serves as our primary reference for robust
statistical methods.

The robust estimators examples in this book suggest that they may
result in better portfolio performance than classical estimators.
But this is not guaranteed -- extensive studies are needed to determine
if and when robust estimators can consistently result in better portfolio
performance. That said, robust estimators serve a very useful purpose
of diagnosing whether or not outliers are influencing a classical
estimator model fit, for, without such comparison, portfolio mangers
will typically be unaware of whether or not outliers are influencing
the classical estimators they typically use. This claim is vividly
demonstrated when comparing robust and classical regression model
fits in Chapters 6 and 7 (time series and cross-section factor models,
respectively).

\section*{Special Focus Topics}

Chapters 10-16 cover a number of special topics that most textbooks
either ignore or address with a light touch. We hope in this book
to rectify these omissions.

In Chapter 10, we cover performance measurement for public and private
funds and the application of statistical process control techniques
(the CUSUM algorithm) to rapidly detect underperformance for actively
managed portfolios. Chapter 11 develops a powerful accounting-based
framework to estimate the expected returns of equity indices and also
addresses the identification of business cycles for purposes of conditional
forecasting. In addition, it covers the estimation of factor risk
premia and touches on the reproducibility and economic viability of
the vast universe of reported anomalies. Some of the material in this
chapter is well known to economists and electrical engineers, particularly
those focused on signal processing, but lacks visibility in finance.
We also explore the links between expected returns, corporate finance
and corporate valuation, all three of which are joined at the hip
intellectually, but are usually treated separately.

Chapter 12 addresses fixed income markets and ends with a discussion
of practical techniques for robust fixed income portfolio management
and immunization of liabilities. Chapter 13 explores the role of options
in portfolios, while Chapter 14 explores end uses of portfolio management
-- in particular, the management of portfolios to meet spending objectives.
More often than not, these topics are treated separately, but it is
impossible to think constructively about one without keeping the other
in mind. Arguably, re-centering one's thoughts around an asset-liability
framework instead of an asset-only framework is the single most important
change that an investor can make, and we hope to nudge our readers
in this direction.

Finally, in Chapters 15 and 16, we return to robust statistics, and
develop ideas presented earlier in the book in greater detail.

\section*{Acknowledgments}

This book, and its accompanying \texttt{PCRA} package, would not have
seen the light of day without help from a vast array of colleagues,
friends and well wishers. Some helped with the development and maintenance
of one or more R packages that support the the book and its associated
reproducibility character. Others read chapters and commented on them,
and yet others granted us formal rights to include datasets in the
\texttt{PCRA} package that are not available to the public. These
contributors are listed in the following paragraphs; we thank them
for their generosity, and apologize in advance to any we might inadvertently
have missed.

First, the people who helped develop and maintain the key R packages
that support this book and its reproducibility features: Xinran Zhao
for \texttt{PortfolioAnalytics}, Brian Peterson for \texttt{PerformanceAnalytics},
Mahmood (Mido) Shamma for \texttt{facmodCS}, Jon Spinney for \texttt{PCRA}
and \texttt{facmodCS}, and Yuhong Liu for \texttt{robustGarch}. Additionally,
Vinav Singh Sancheti wrote the code for the CUSUM algorithm.

Next, our readers: David Carino, Roland Lochoff, George Mylnikov,
Dale Rosenthal, Nikolay Ryabkov, Jon Spinney, Stoyan Stoyanov and
Jarrod Wilcox read and commented on multiple chapters. Their inputs
proved invaluable and led us to substantially rewrite every single
chapter they read. Stoyan, in particluar, is a long-time friend and
LyX/\LaTeX{}\ whiz who helped clarify our thinking, and without
whose guidance we would have struggled to get the book ready for production.
Doug and Tom had the benefit of outstanding Ph. D students and teaching
assistants (Yindeng Jiang, Shengyu Zhang, Ercument Cahan, Cheng Xu
and Yuchen Han) who corrected their errors with grace. It's fair to
say that we learned as much from our readers, students and teaching
assistants as they did from us.

Third, the wonderful people who provided us tremendous assistance
in obtaining various rights to use data from their companies in this
book and its supporting \texttt{PCRA} package: Taras Zlupko, Janet
Eder, Colin O'Connell, Diane Roman, and Jeffrey Ostrom (all of \CRSPTM),
Joel Nadelman, Kirk Wang, Howard Silverblatt, Guillermo Ruiz-Rico,
and Jonathan Campbell (all of S\&P, and SPGMI), Mike Cassidy and Elizabeth
Weinstein (ICE), Mallory Denney and Namita Kapaley (LSE Group), Geralyn
Endo (CBOE), Lisa Schirf and Xinwei Wang (Tradeweb), all helped us
obtain data that is provided free for non--commercial use in the
\texttt{PCRA} package. Their generosity contributed substantially
to our efforts in creating this book.

And finally, our families. Our parents, spouses, children and grandchildren
(with intra-author variations!) put up with our absences for long
stretches of time while we worked on this book. Without their unconditional
love and support, this book would not have seen the light of day.
One of us (TP) adds a very special thank you to the two strangers
who pulled him out of the clutches of a rip tide off the coast of
Montauk in September 2009. Were it not for their bravery, it would
have been someone else who wrote certain chapters.

\section*{Disclaimers Related to PCRA Datasets}

Thanks to the generosity of a number of industry sponsors, we have
been able to supply historical and market data for a variety of securities,
indices and factors. In all cases, the data provided with the \texttt{PCRA}
package  that accompanies this book is \emph{strictly for educational
(i.e. non-commercial) use} -- anyone who wishes to use the data for
commercial purposes must enter into an appropriate licensing agreement
with the relevant vendor. Furthermore, the data cannot be redistributed
or published elsewhere. Please read the relevant vendor disclosures
below.

\subsection*{CRSP®}

The University of Chicago, the Booth School of Business, CRSP ® ,
and its third party suppliers make no warranty of any kind, express
or implied, with respect to the merchantability, fitness, condition,
use or appropriateness for permittee’s purposes of the CRSP® data,
or any other purpose, and all CRSP® data is supplied on an “as is”
basis. Use of CRSP® data is subject to the disclaimers provided at
\url{https://www.crsp.org/policies-statements/}.

\subsection*{SPGMI}

Standard and Poors Global Market Intelligence (SPGMI) is a business
unit of Standard and Poors Global that provides multi-asset class
and real-time data, research, news and analytics. Use of SPGMI data
is subject to the disclaimers provided at \url{https://www.spglobal.com/market-intelligence/en/legal/disclosures}.

\subsection*{LSE / FTSE Russell}

FTSE Russell is a trading name of certain of the LSE Group companies.
“Russell®”/“FTSE Russell®”/“Mergent®”, are/is trade mark(s) of the
relevant LSE Group companies and is/are used by any other LSE Group
company under license. All rights in the FTSE Russell indexes or data
vest in the relevant LSE Group company which owns the index or the
data. Neither LSE Group nor its licensors accept any liability for
any errors or omissions in the indexes or data and no party may rely
on any indexes or data contained in this communication. No further
distribution of data from the LSE Group is permitted without the relevant
LSE Group company’s express written consent. The LSE Group does not
promote, sponsor or endorse the content of this communication.

\subsection*{ICE}

Source: ICE Data Indices, LLC (\textquotedbl ICE Data\textquotedbl ),
used with permission. THE ICE DATA INDICES ARE PROVIDED BY ICE DATA
FOR EDUCATIONAL PURPOSES. ICE DATA, ITS AFFILIATES AND THEIR RESPECTIVE
THIRD PARTY SUPPLIERS DISCLAIM ANY AND ALL WARRANTIES AND REPRESENTATIONS,
EXPRESS AND/OR IMPLIED, INCLUDING ANY WARRANTIES OF MERCHANTABILITY
OR FITNESS FOR A PARTICULAR PURPOSE OR USE, INCLUDING THE INDICES,
INDEX DATA AND ANY DATA INCLUDED IN, RELATED TO, OR DERIVED THEREFROM.
NEITHER ICE DATA, ITS AFFILIATES NOR THEIR RESPECTIVE THIRD PARTY
SUPPLIERS SHALL BE SUBJECT TO ANY DAMAGES OR LIABILITY WITH RESPECT
TO THE ADEQUACY, ACCURACY, TIMELINESS OR COMPLETENESS OF THE INDICES
OR THE INDEX DATA OR ANY COMPONENT THEREOF, AND THE INDICES AND INDEX
DATA AND ALL COMPONENTS THEREOF ARE PROVIDED ON AN \textquotedbl AS
IS\textquotedbl{} BASIS AND YOUR USE IS AT YOUR OWN RISK. ICE DATA,
ITS AFFILIATES AND THEIR RESPECTIVE THIRD PARTY SUPPLIERS DO NOT SPONSOR,
ENDORSE, OR RECOMMEND ANY TEXTBOOKS, PRODUCTS OR SERVICES OF SPRINGER
PUBLISHERS, R. DOUGLAS MARTIN, THOMAS K. PHILIPS, BERND SCHERER, OR
KIRK LI.

\subsection*{Tradeweb}

Tradeweb OTR--OFTR Spread Data is provided by Tradeweb Markets LLC
for educational purposes. only. Tradeweb Markets LLC makes no representation
or warranty of any kind, express or implied, including regarding the
accuracy, adequacy, validity, reliability, availability or completeness
of the Tradeweb OTR--OFTR Spread Data. All rights reserved.

\subsection*{S\&P}

S\&P® and S\&P 500® are registered trademarks of Standard \& Poor's
Financial Services LLC, and Dow Jones® is a registered trademark of
Dow Jones Trademark Holdings LLC. © 2020 S\&P Dow Jones Indices LLC,
its affiliates and/or its licensors. All rights reserved.

Reproduction of any information, data or material, including ratings
(\textquotedbl Content\textquotedbl ) in any form is prohibited
except with the prior written permission of the relevant party. Such
party, its affiliates and suppliers (\textquotedbl Content Providers\textquotedbl )
do not guarantee the accuracy, adequacy, completeness, timeliness
or availability of any Content and are not responsible for any errors
or omissions (negligent or otherwise), regardless of the cause, or
for the results obtained from the use of such Content. In no event
shall Content Providers be liable for any damages, costs, expenses,
legal fees, or losses (including lost income or lost profit and opportunity
costs) in connection with any use of the Content. A reference to a
particular investment or security, a rating or any observation concerning
an investment that is part of the Content is not a recommendation
to buy, sell or hold such investment or security, does not address
the suitability of an investment or security and should not be relied
on as investment advice. Credit ratings are statements of opinions
and are not statements of fact.

\subsection*{CBOE LIVEVOL}

CBOE LIVEVOL data is provided \textquotedbl as is\textquotedbl{}
without warranty of any kind, either express or implied, including,
without limitation, any warranty with respect to accuracy, completeness,
timeliness, noninfringement, merchantability, or fitness for a particular
purpose. Neither LIVEVOL, nor any provider of data to LIVEVOL, nor
any of their respective affiliates, nor their respective directors,
officers, employees, contractors, and agents shall have any liability
of any kind (including, but not limited to, for any direct, indirect,
incidental, special, consequential, or punitive damages or any damages
for lost profits or lost opportunities and whether based upon contract,
tort, warranty, or otherwise) for any inaccuracies, omissions, human
or machine errors, or other irregularities in the data or for any
cessation, discontinuance, failure, malfunction, delay, suspension,
interruption, or termination of, or with respect to, the provision
of the data to subscriber. The data is not, and should not be construed
as financial, legal or other advice of any kind, nor should it be
regarded as an offer or as a solicitation of an offer to buy, sell
or otherwise deal in any investment.

\printbibliography
\end{document}
